\documentclass{beamer}
\usetheme{Madrid}
\usecolortheme{default}

\title[ESP32 Linux Compatibility Layer]{Architecting a POSIX-Compliant Runtime Environment on ESP32}
\subtitle{A Unikernel Approach to Dynamic ELF Execution}
\author{Dhruv}
\institute[University]{Department of Computer Science}
\date{\today}

\begin{document}

% Title Slide
\begin{frame}
    \titlepage
\end{frame}

% Outline
\begin{frame}{Outline}
    \tableofcontents
\end{frame}

% Section 1: The Problem
\section{The Problem}
\begin{frame}{The Problem: Static vs. Dynamic}
    \begin{itemize}
        \item \textbf{Embedded Systems (RTOS):}
        \begin{itemize}
            \item Highly efficient, deterministic.
            \item \textbf{Static:} Monolithic firmware. Adding a feature requires recompiling and reflashing the entire OS.
            \item "OTA Updates" replace the whole image.
        \end{itemize}
        \vspace{0.5cm}
        \item \textbf{General Purpose OS (Linux):}
        \begin{itemize}
            \item Dynamic. Applications are separate binaries.
            \item Can load/unload code at runtime (`exec`).
            \item Requires \textbf{MMU} (Memory Management Unit) and significant RAM.
        \end{itemize}
        \vspace{0.5cm}
        \item \textbf{The Gap:} Modern IoT (Edge Computing, C2, Drones) needs the dynamism of Linux on the hardware constraints of an MCU.
    \end{itemize}
\end{frame}

% Section 2: The Solution
\section{The Solution}
\begin{frame}{The Solution: A Thin Linux Compatibility Layer}
    \begin{block}{Architecture: The Unikernel Model}
        Instead of running a full kernel, we implement a \textbf{Library OS} that translates Linux System Calls into ESP-IDF primitives.
    \end{block}

    \begin{itemize}
        \item \textbf{No-MMU Constraint:} We cannot implement `fork()` or virtual memory.
        \item \textbf{Single Address Space:} The Kernel and Guest App share the same physical RAM.
        \item \textbf{ELF Loader:} A custom loader parses standard `.elf` binaries, allocates IRAM (Instruction RAM) for code, and performs \textit{Xtensa-specific relocations} at load time.
    \end{itemize}
\end{frame}

\begin{frame}{System Architecture}
    \begin{center}
    % Conceptual Diagram
    \texttt{
    [ Guest ELF Application ] (printf, open, socket)\\\n             |\\
             v\\
    [ Syscall Shim Layer ] (Path Translation, Errno)\\\n             |\\
             v\\
    [ ESP-IDF / FreeRTOS ] (LwIP, VFS, Drivers)\\\n             |\\
             v\\
    [ ESP32 Hardware ]
    }
    \end{center}
\end{frame}

% Section 3: Demo 1
\section{Demo 1: Distributed C2 Botnet}
\begin{frame}{Demo 1: Distributed Command & Control}
    \begin{block}{Scenario}
        An "Attacker" sends a compiled binary payload to an edge node (ESP32) over the network. The node executes it immediately and streams output back.
    \end{block}

    \begin{columns}
        \column{0.5\textwidth}
        \textbf{Workflow:}
        \begin{enumerate}
            \item Master sends ELF Header + Data via TCP.
            \item ESP32 saves to `/linux/payload.elf`.
            \item \textbf{Redirection Shim:} `dup2(socket, stdout)` allows `printf` to stream over WiFi.
            \item `execve()` spawns a new task.
        \end{enumerate}

        \column{0.5\textwidth}
        \textbf{Key Tech:}
        \begin{itemize}
            \item \textbf{shim\_write:} Intercepts FD 1 (stdout) and dual-routes data to UART and TCP.
            \item \textbf{Spawn Model:} Simulates `fork+exec` using `xTaskCreate`.
        \end{itemize}
    \end{columns}
\end{frame}

% Section 4: Demo 2
\section{Demo 2: Collision Avoidance}
\begin{frame}{Demo 2: Cyber-Physical System (V2X)}
    \begin{block}{Scenario}
        An Edge-OS Collision Avoidance System. A Linux app performs real-time vector math to detect crashes and triggers hardware alarms.
    \end{block}

    \begin{itemize}
        \item \textbf{Telemetry Bridge:} Web Dashboard sends JSON (GPS, Speed) $\rightarrow$ ESP32 $\rightarrow$ UDP $\rightarrow$ Linux App.
        \item \textbf{Vector Math (User Space):}
        $$ \vec{V}_{rel} = \vec{V}_{target} - \vec{V}_{ego} $$
        $$ TTC = \frac{|\vec{P}_{target} - \vec{P}_{ego}|}{|\vec{V}_{rel}|} $$
        \item \textbf{Hardware Abstraction (Kernel Space):}
        \begin{itemize}
            \item App writes to `/dev/buzzer`.
            \item Custom VFS Driver translates `write()` to `ledc_set_duty()` (PWM).
        \end{itemize}
    \end{itemize}
\end{frame}

% Conclusion
\section{Conclusion}
\begin{frame}{Conclusion & Future Work}
    \begin{itemize}
        \item \textbf{Achievement:} Successfully implemented a POSIX-compliant runtime on ESP32 capable of running standard C applications.
        \item \textbf{Impact:} Enables "Write Once, Run Anywhere" for IoT edge logic without firmware reflashing.
        \item \textbf{Future Work:}
        \begin{itemize}
            \item Implement `setjmp`/`longjmp` for fuller signal handling.
            \item Add support for dynamic shared libraries (`.so`).
            \item Port a lightweight scripting language (Lua/MicroPython) as a guest app.
        \end{itemize}
    \end{itemize}
\end{frame}

\end{document}
